\xiaosi

本文以神经电路策略（NCP）为核心，围绕"可解释自动驾驶规划"展开研究，主要内容包括以下四个部分：

\textbf{（1）NCP规划头的性能验证。}将NCP集成为Double DQN的Q-network，针对NCP的ODE约束问题提出两层Q-head方案（用全部神经元隐状态经MLP映射Q-values），在Highway-env的4个场景（高速巡航、匝道合流、交叉路口、环岛）上与Random、MLP、LSTM、GRU、全连接LTC共6种模型进行系统对比。采用3个随机种子确保统计显著性，评估维度包括累计奖励、碰撞率、参数量、推理延迟和参数效率。

\textbf{（2）NCP布线拓扑的进化搜索。}设计结构化搜索空间（6个离散基因：神经元数量、扇出/扇入度、循环连接数），使用进化策略（锦标赛选择+均匀交叉+高斯变异）进行搜索。采用安全感知的多目标适应度函数，在多场景联合评估下自动发现兼顾性能与安全性的最优布线方案。

\textbf{（3）消融实验。}包括三组消融：Q-head消融（验证两层Q-head对NCP性能的必要性）、ODE展开数消融（确定精度-效率最优权衡点）、网络规模消融（探究NCP性能与规模的关系）。

\textbf{（4）多层次可解释性分析。}设计涵盖5个层次的分析框架：结构层（拓扑统计量对比）、动态层（command神经元激活与驾驶决策对齐）、安全层（碰撞前关键时刻异常激活识别）、因果层（逐神经元mask的功能归因矩阵）、泛化层（跨场景激活降维对比）。

本文的技术路线为：首先在Highway-env上搭建统一的实验框架，验证NCP的基线性能；然后通过进化搜索优化布线拓扑；接着通过消融实验分析各设计要素的贡献；最后通过多层次可解释性分析展示NCP的独特价值。

本文的预期结果包括：（1）NCP以不到1万参数达到与LSTM、GRU等传统循环网络相当甚至更优的决策性能，验证其作为轻量规划头的可行性；（2）进化搜索发现的最优布线拓扑在累计奖励和碰撞率上优于手工设计的默认拓扑；（3）消融实验定量揭示两层Q-head、ODE展开数和网络规模对NCP性能的贡献程度；（4）多层次可解释性分析表明NCP的结构化稀疏连接能够产生与驾驶语义对齐的神经元激活模式，为自动驾驶规划系统的安全合规提供结构性支撑。
